\begin{abstract}
Identifying the hadronic decays of highly boosted top quarks --- ``top
tagging'' --- is a benchmark problem in collider physics and a proving
ground for modern machine learning. This thesis compares four classifiers
on the public Top-Quark Tagging Reference Dataset of 14\,TeV
proton--proton collisions: a gradient-boosted decision tree (BDT) on
engineered jet-substructure variables, a convolutional neural network (CNN)
on jet images, the ParticleNet graph network, and the Particle Transformer.
Beyond reproducing the usual ranking, the study asks two questions of direct
practical value: how much training data each architecture needs (data
efficiency), and whether its output scores can be trusted as probabilities
(calibration). Using a reproducible pipeline trained on a free cloud GPU, we
sweep the training set from $5\times10^{3}$ to $5\times10^{4}$ jets and
measure test AUC, background rejection at fixed signal efficiency, the
Expected Calibration Error (ECE), and the effect of temperature scaling. The
Particle Transformer attains the best discrimination (AUC $=0.978$,
background rejection $\approx136$ at $50\%$ signal efficiency), followed by
ParticleNet, the CNN and the BDT, reproducing the published ordering. The
data-efficiency study reveals a clear crossover: the BDT is strongest at the
smallest training size but is overtaken by the deep models as data grows,
the Particle Transformer rising from weakest to strongest. The calibration
study shows that ranking and calibration are decoupled --- the jet-image CNN
is the worst-calibrated model (ECE $=0.078$) despite its mid-pack ranking,
whereas the top-ranked transformer is comparatively well-calibrated; a
single-parameter temperature rescaling corrects all models except the CNN.
The complete pipeline, trained models and documents are released as open
source.

\vspace{2mm}
\noindent\textbf{Keywords:} top tagging, jet substructure, deep learning,
ParticleNet, Particle Transformer, data efficiency, probability calibration.
\end{abstract}
